
\noindent \uline{Lorsqu'on considère les deux membres de} (7.3.1) \uline{comme des modules sur} \(K^{\bullet}(X)\) \uline{grâce aux homomorphismes d'anneaux} \(K^{\bullet}(X)\xrightarrow{j^*}K^{\bullet}(X_t)\) \uline{et} \(K^{\bullet}(X)\xrightarrow{i^*}K^{\bullet}(X_s)\) (IV 2.7 b) \uline{et aux structures de modules de} \(K_\bullet(X_t)\) \uline{et} \(K_\bullet(X_s)\) \uline{sur} \(K^{\bullet}(X_t)\) \uline{et} \(K^{\bullet}(X_s)\) \uline{respectivement} (IV 2.10), \uline{l'homomorphisme} \(\sigma_\bullet\) \uline{est un homomorphisme de} \(K^{\bullet}(X)\)-\uline{modules}.

La dernière assertion provient du fait que \(i^*:K_\bullet(X)\to K_\bullet(X_s)\) est un homomorphisme de \(K^{\bullet}(X)\)-modules, ainsi que l'épimorphisme \(j^*:K_\bullet(X)\longrightarrow K_\bullet(X_t)\), en vertu de (IV 2.12.2).

Nous allons passer en revue d'autres propriétés du morphisme de spécialisation.

\subsection*{7.4. Compatibilité aux filtrations.}
Munissons les \(K_\bullet\) de la filtration \(\operatorname{Fil}_\bullet K_\bullet\) par la dimension des supports (1.1.1). Alors, si \(f\) est de type fini, plus généralement si \(X\) est caténaire, \emph{l'homomorphisme de spécialisation est compatible avec ces filtrations}, et induit par suite un homomorphisme (appelé également homomorphisme de spécialisation) :
\begin{equation}\tag{7.4.1}
  \sigma_\bullet:\Gr_\bullet(X_t)\longrightarrow \Gr_\bullet(X_s).
\end{equation}

Pour le voir, on est ramené par 1.14 à prouver que si \(Y_t\) est un sous-schéma intègre de dimension \(d\) de \(X_t\), alors \(\sigma_\bullet(\cl_\bullet(Y_t))\) est de filtration \(\leq d\). Or soit \(Y\) d'adhérence schématique de \(Y_t\) dans \(X\), de sorte que \(Y\) est plat sur \(S\) (EGA IV 2.8.1). Il en résulte aisément que \(i^*(\cl_\bullet(Y))=\cl_\bullet(Y_s)\), où les deux \(\cl_\bullet\) sont pris respectivement dans \(K_\bullet(X)\) et dans \(K_\bullet(X_s)\), d'où \(\sigma_\bullet(\cl_\bullet(Y_t))=\cl_\bullet(Y_s)\). On vérifie aisément (grâce au fait que \(X\) est caténaire, donc \(Y\) caténaire) que \(\dim Y_s\leq \dim Y_t\) ; dans le cas où \(f\) est de type fini, cela résulte aussitôt, par réduction au cas \(X\) intègre, de EGA IV 7.1.13. D'où la conclusion annoncée.

